% kotlin-coroutines-internals.tex — how coroutines work underneath: the CPS transform
% (Continuation param, Any? return, COROUTINE_SUSPENDED), the generated state machine (label,
% spilled locals), ContinuationInterceptor + dispatch, CoroutineContext as an indexed set, Job
% states and the parent/child tree, structured-concurrency rules, cooperative cancellation (and
% the runCatching trap), exception propagation launch vs async, coroutineScope vs supervisorScope,
% CoroutineExceptionHandler, NonCancellable, suspendCancellableCoroutine, dispatchers +
% limitedParallelism. Diagram: state machine of a 2-suspension function + the Job tree.
% Goes deeper than android-coroutines-platform (usage level) — does not repeat it.
% Sources (checked 2026-09-26):
%   https://github.com/Kotlin/KEEP/blob/master/proposals/coroutines.md  (CPS, state machines)
%   https://kotlinlang.org/docs/exception-handling.html · https://kotlinlang.org/docs/cancellation-and-timeouts.html
%   https://kotlinlang.org/api/kotlinx.coroutines/kotlinx-coroutines-core/kotlinx.coroutines/-job/  (states table)
%   https://kotlinlang.org/api/kotlinx.coroutines/kotlinx-coroutines-core/kotlinx.coroutines/-dispatchers/-i-o.html
%   https://kotlinlang.org/docs/coroutine-context-and-dispatchers.html
% @labels: area=cross-platform kind=concept platform=android topic=concurrency,language discussion=117
% @tags: kotlin-coroutines, continuation-passing-style, state-machine, coroutine-context, job-hierarchy, structured-concurrency, cooperative-cancellation, exception-propagation, supervisorscope, dispatchers, suspendcancellablecoroutine
\documentclass[8pt]{extarticle}
\usepackage{printup-sheet}
\usepackage{array}
\usepackage{tabularx}

% Kotlin listings language; literate = one \hbox per char so 0xProto ligatures never form
\lstdefinelanguage{KotlinSheet}{
  morekeywords={fun,val,var,class,object,data,sealed,interface,suspend,when,is,in,out,
    override,private,return,if,else,try,catch,finally,throw,while,null,this,it,inline},
  sensitive=true, morecomment=[l]{//}, morecomment=[s]{/*}{*/}, morestring=[b]",
  literate={->}{{\hbox{-}\hbox{>}}}2 {?.}{{\hbox{?}\hbox{.}}}2 {?:}{{\hbox{?}\hbox{:}}}2
           {!!}{{\hbox{!}\hbox{!}}}2 {==}{{\hbox{=}\hbox{=}}}2 {!=}{{\hbox{!}\hbox{=}}}2
           {::}{{\hbox{:}\hbox{:}}}2 {&&}{{\hbox{\&}\hbox{\&}}}2}
\lstset{basicstyle=\ttfamily\scriptsize, aboveskip=2pt, belowskip=2pt}

\newcommand\ar{\texttt{\hbox{-}\hbox{>}}}

\tikzset{
  st/.style={box, font=\tiny, inner sep=1.2pt, minimum height=4.2mm, text width=#1, align=left},
  st/.default=24mm,
  jb/.style={box, font=\tiny, inner sep=1.2pt, minimum height=4mm, minimum width=13mm},
  ok/.style={jb, draw=sheetGreen, fill=sheetGreen!12},
  ko/.style={jb, draw=sheetRed, fill=sheetRed!10},
  gone/.style={jb, draw=sheetGrey, fill=sheetGrey!10, dashed},
  lbl/.style={font=\tiny, text=black!75, inner sep=1pt, align=center},
  hd/.style={font=\bfseries\scriptsize, anchor=west},
}
\newcolumntype{L}{>{\raggedright\arraybackslash}X}

\begin{document}

\sheettitle{Kotlin coroutines — internals}{cross-platform · folio}

\oneliner{The compiler rewrites every \texttt{suspend fun} in \textbf{continuation-passing
style}: it gains a hidden \texttt{Continuation} parameter, returns \texttt{Any?}, and its body
becomes a \textbf{state machine} whose \texttt{label} says where to resume. A suspension
returns \texttt{COROUTINE\_SUSPENDED} all the way up and frees the thread; later something calls
\texttt{resumeWith}, the \textbf{dispatcher} picks a thread, the machine jumps to the next
label. \texttt{Job}, scopes, cancellation = \emph{library} (kotlinx.coroutines) on top.}

\vspace{2pt}
\noindent\begin{tikzpicture}[sheet]
  % ---------- A: state machine ----------
  \node[hd] at (-0.3,3.35) {\texttt{suspend fun load(id)\{ val u = user(id); val p = posts(u); return Ui(u,p) \}}};
  \node[st=27mm, draw=sheetBrown, fill=sheetBrown!8] (obj) at (1.25,2.7)
    {\textbf{\texttt{load\$1 : ContinuationImpl}}\\\texttt{label}, \texttt{result},
     \texttt{L\$0} (spilled \texttt{u}), \texttt{completion}};
  \node[st=19mm] (s0) at (0.75,1.55) {\textbf{label 0}: \texttt{label=1};\\call \texttt{user(id,}\\\texttt{this)}};
  \node[st=19mm] (s1) at (3.95,1.55) {\textbf{label 1}: \texttt{u=result};\\\texttt{L\$0=u; label=2};\\call \texttt{posts(u,this)}};
  \node[st=19mm, draw=sheetGreen, fill=sheetGreen!10] (s2) at (7.4,1.55) {\textbf{label 2}: \texttt{u=L\$0};\\\texttt{p=result};\\\texttt{return Ui(u,p)}};
  \draw[hot] (s0) -- node[lbl, above]{\texttt{resumeWith}} node[lbl, below]{\texttt{(u)}} (s1);
  \draw[hot] (s1) -- node[lbl, above]{\texttt{resumeWith}} node[lbl, below]{\texttt{(p)}} (s2);
  \node[lbl, text=sheetRed] (su0) at (0.75,0.6) {returns\\\texttt{COROUTINE\_SUSPENDED}};
  \node[lbl, text=sheetRed] (su1) at (3.95,0.6) {returns\\\texttt{COROUTINE\_SUSPENDED}};
  \draw[->, thick, dashed, sheetRed] (s0) -- (su0);
  \draw[->, thick, dashed, sheetRed] (s1) -- (su1);
  \node[lbl, text=sheetGreen!60!black] (dn) at (7.4,0.6) {\texttt{completion.resumeWith(}\\\texttt{Result.success(ui))}};
  \draw[flow, draw=sheetGreen] (s2) -- (dn);
  \draw[flow, dashed, draw=sheetBrown] (obj.east) -- ++(1.2,0) node[lbl, right, text=sheetBrown, align=left]
    {one object per call, reused for every resume\\\texttt{invokeSuspend(result)} = \texttt{switch(label)}};
  \node[lbl, text=sheetBlue, align=left, anchor=west] at (-0.3,0.0) {resume: callback thread $\to$ \texttt{resumeWith} $\to$
    \texttt{DispatchedContinuation} $\to$ \texttt{dispatcher.dispatch()}\\$\to$ \texttt{invokeSuspend} on the chosen thread.
    \textcolor{sheetGrey}{Fast path: callee returned a value, not the marker $\to$ fall through.}};
  \draw[sheetGrey!40] (8.75,3.3) -- (8.75,-0.45);
  % ---------- B: Job tree ----------
  \node[hd] at (8.9,3.35) {Job tree — failure up, cancellation down};
  \node[jb, draw=sheetBlue, minimum width=24mm] (root) at (12.6,2.85) {\texttt{coroutineScope \{\}} (Job)};
  \node[ko] (p) at (10.6,2.05) {launch A: \textbf{Cancelling}};
  \node[ok, draw=sheetBlue, fill=sheetBlue!8] (sv) at (14.7,2.05) {\texttt{supervisorScope}};
  \node[ko] (c1) at (9.7,1.2) {async a1 \textbf{throws IOE}};
  \node[gone] (c2) at (11.55,1.2) {launch a2};
  \node[ko] (d1) at (13.8,1.2) {launch s1 throws};
  \node[ok] (d2) at (15.7,1.2) {launch s2 lives};
  \draw[flow] (root) -- (p); \draw[flow] (root) -- (sv);
  \draw[flow] (p) -- (c1); \draw[flow] (p) -- (c2);
  \draw[flow] (sv) -- (d1); \draw[flow] (sv) -- (d2);
  \draw[hot, draw=sheetRed] (c1.north) to[bend left=35] node[lbl, left, text=sheetRed]{1 fail} (p.west);
  \draw[hot, draw=sheetRed] (p.north) to[bend left=25] node[lbl, above left, text=sheetRed]{2 fail} (root.west);
  \draw[hot, draw=sheetOrange] (p.south east) -- node[lbl, right, text=sheetOrange]{3 CE} (c2.north);
  \draw[hot, draw=sheetOrange, dashed] (root.south east) to[bend left=15] node[lbl, right, text=sheetOrange]{4 cancel} (sv.north);
  \node[lbl, text=sheetRed, align=left, anchor=west] at (13.1,0.6) {s1's error stops at the supervisor:\\goes to \emph{s1's} \texttt{CoroutineExceptionHandler}};
  \node[lbl, text=sheetBrown, align=left, anchor=west] at (8.9,0.6) {5 scope waits for all children,\\then \textbf{rethrows IOE} to caller};
  \node[lbl, align=left, anchor=west] at (8.9,0.0) {States: \texttt{New}\,(LAZY) $\to$ \texttt{Active} $\to$ \texttt{Completing}
    (waits children) $\to$ \texttt{Completed};\\cancel/fail $\to$ \texttt{Cancelling} $\to$ \texttt{Cancelled}.
    \texttt{isActive}: Active/Completing only.};
\end{tikzpicture}

\begin{multicols}{2}
\raggedright

\section{The CPS transform}
\begin{itemize}\raggedright
  \item \texttt{suspend fun f(x: A): T} $\to$ JVM \texttt{Object f(A x, Continuation<? super T> c)}.
        Return is \texttt{T} \emph{or} the marker \texttt{COROUTINE\_SUSPENDED}
        (\texttt{kotlin.coroutines.intrinsics}) — hence \texttt{Any?}.
  \item \texttt{interface Continuation<in T> \{ val context: CoroutineContext;
        fun resumeWith(result: Result<T>) \}} — a typed callback carrying its context.
  \item Only functions that \emph{call} suspend functions get a state machine; one
        \texttt{ContinuationImpl} subclass per function, \textbf{locals live across a
        suspension are spilled} into fields (\texttt{L\$0}, \texttt{I\$0}…\unverified), the rest stay on
        the stack. A suspended coroutine = that heap object, \textbf{no thread held}.
\end{itemize}

\section{Context \& dispatch}
\begin{itemize}\raggedright
  \item \texttt{CoroutineContext} = \textbf{indexed set} of \texttt{Element}s keyed by
        \texttt{Key}: \texttt{Job}, \texttt{ContinuationInterceptor} (the dispatcher),
        \texttt{CoroutineName}, \texttt{CoroutineExceptionHandler}. \texttt{a + b}: the right
        side wins per key; \texttt{ctx[Job]}, \texttt{minusKey}.
  \item A child's context = parent context + builder arg + a \textbf{new child Job}.
        Passing \texttt{Job()} as an argument \textbf{breaks the parent link}.
  \item Dispatcher = interceptor: \texttt{interceptContinuation} wraps so \texttt{resume}
        $\to$ \texttt{dispatch()}.
  \item \texttt{Default}: \#cores (min 2), CPU work. \texttt{IO}: 64 or \#cores, whichever is
        larger — \textbf{shares threads with Default}, so \texttt{withContext(IO)} from Default
        often does not switch thread. \texttt{IO.limitedParallelism(n)} is \emph{elastic}
        (beyond the 64); \texttt{Default.limitedParallelism(n)} is a view within Default.
        \texttt{Main.immediate}: no re-post if already on main.
\end{itemize}

\section{Example — bridge a callback, cancel cleanly}
\begin{lstlisting}[language=KotlinSheet]
suspend fun Call.await(): Response =
  suspendCancellableCoroutine { cont ->
    cont.invokeOnCancellation { cancel() }   // propagate inward
    enqueue(object : Callback {
      override fun onResponse(c: Call, r: Response) = cont.resume(r)
      override fun onFailure(c: Call, e: IOException) =
        cont.resumeWithException(e)
    })
  }
suspend fun crunch(xs: List<Int>) = withContext(Dispatchers.Default) {
  for (x in xs) { ensureActive(); heavy(x) }  // cooperative
}
try { work() } finally {
  withContext(NonCancellable) { db.flush() } // may suspend
}
runCatching { fetch() }          // TRAP: also catches CancellationException
  .onFailure { if (it is CancellationException) throw it }
\end{lstlisting}

\columnbreak

\section{Structured concurrency \& cancellation}
\begin{itemize}\raggedright
  \item Rules: every coroutine has a parent Job; a parent \textbf{completes only after its
        children}; \textbf{cancelling a parent cancels all children}; a child failing with a
        non-\texttt{CancellationException} \textbf{cancels its parent} and so its siblings.
  \item \texttt{job.cancel()} takes only a \texttt{CancellationException} = \emph{normal}:
        parent \textbf{not} cancelled, handlers ignore it.
  \item Cancellation is \textbf{cooperative}: a flag that bites at the next
        \emph{cancellable} suspension point (\texttt{delay}, \texttt{yield},
        \texttt{withContext}, \texttt{await}…) by throwing \texttt{CancellationException}.
        CPU loops check \texttt{isActive}/\texttt{ensureActive()}/\texttt{yield()}.
  \item Suspending in \texttt{finally} after cancel throws again $\to$
        \texttt{withContext(NonCancellable)}. \texttt{withTimeout} throws a CE subclass.
\end{itemize}

\section{Exceptions — who sees what}
{\scriptsize
\noindent\begin{tabularx}{\linewidth}{@{}>{\raggedright\arraybackslash}p{22mm}L@{}}
\toprule
\textbf{Where} & \textbf{What happens to an exception} \\
\midrule
root \texttt{launch} & uncaught $\to$ its \texttt{CoroutineExceptionHandler}, else the thread's
  uncaught handler (Android: crash) \\
root \texttt{async} & stored in the \texttt{Deferred}; thrown at \texttt{await()};
  a CEH on it has \textbf{no effect} \\
any child & delegated to the parent, up to the root — a CEH on a child is \textbf{never used};
  an un-awaited child \texttt{async} still fails the parent \\
\texttt{coroutineScope} & cancels the rest, waits, \textbf{rethrows} to the caller —
  a normal \texttt{try/catch} around it works \\
\texttt{supervisorScope} / \texttt{SupervisorJob} & child failure is not propagated; direct
  children use \emph{their own} CEH (like roots) \\
several fail & first exception wins, the rest are \texttt{suppressed} \\
\bottomrule
\end{tabularx}}

\section{Traps}
\begin{itemize}\raggedright
  \trap{\texttt{runCatching}/\texttt{catch (e: Exception)} swallows
        \texttt{CancellationException}: the coroutine runs on after cancel.}
  \trap{\texttt{launch(SupervisorJob())} in a scope: that Job becomes the parent — nothing
        supervised, scope cancel no longer reaches it.}
  \trap{\texttt{try/catch} around \texttt{launch \{\}} catches nothing — the exception goes up
        the Job tree, not up the call stack.}
  \trap{Blocking calls ignore cancel — use \texttt{runInterruptible}.}
\end{itemize}

\textbf{Remember:} \emph{suspend = callback + state machine · failure up, cancel down ·
CE is normal · handler at the root.}


\end{multicols}

\noindent{\footnotesize\color{sheetGrey}\textit{Related:} android-coroutines-platform ·
kotlin-flow-deep · kotlin-under-the-hood · kmp-stack-concurrency · kotlin-language ·
concurrency (Swift) · swift6-strict-concurrency}

\end{document}
